\section{Instantiation}\label{sec:agg-creds-instantiation}

\paragraph{Notation.} We use cyclic groups $\group{1}$, $\group{2}$, and
$\group{T}$ of prime order $p$.  The groups are equipped with an efficiently
computable pairing written as $e \colon \group{1} \times \group{2} \rightarrow \group{T}$.
For all $a \in \group{1}$, $\widetilde{b} \in \group{2}$, and $x,y \in \Zp$,
bilinearity means that the following holds:
    %
\begin{equation*}
     %
  e(a^x,\widetilde{b}^{,y}) = e(a,\widetilde{b})^{xy}.
    %
\end{equation*}
     %
The pairing is non-degenerate.  For generators $g \in \group{1}$ and
$\widetilde{g} \in \group{2}$, we have that the element $e(g,\widetilde{g})$
generates $\group{T}$. We work in the asymmetric pairing setting, and so do not
assume the existence of an efficiently computable homomorphism between
$\group{1}$ and $\group{2}$.

Elements of $\group{1}$ are written using lower case letters.  Elements of
$\group{2}$ are distinguished using a tilde, while elements of $\Zp$ are
written using lower case Greek letters or otherwise specified explicitly.  We
use generators $g \in \group{1}$ and $\widetilde{g} \in \group{2}$.

We also use cryptographic hash functions 
 $ \hashone \colon {0,1}^* \rightarrow \group{1}$ and $\hashtwo
  \colon {0,1}^* \rightarrow \group{1}$.

Bilinearity allows relations between exponents in the source groups to be
checked through equalities in $\group{T}$.  For example, an equality of the
form
    %
\begin{equation*}
     %
  e(a,\widetilde{g}) = e(g,\widetilde{b})
    %
\end{equation*}
     %
shows that $a$ and $\widetilde{b}$ contain the same exponent relative to their
respective generators, without requiring that exponent to be recovered.  Our
constructions use such pairing equations to verify signatures and other
algebraic relations across $\group{1}$ and $\group{2}$.

\subsection{Our tag based aggregate signature scheme ($\ATS$)}
\label{sec:tag-scheme}

Our tag based aggregate signature scheme is an enhanced aggregate BLS signature
scheme that signs pairs of messages and tags~\cite{JC:BonLynSha04,EC:BGLS03}.
The key differences are tag restricted aggregation and ability to randomise
signatures.  Our scheme is comprised of the following algorithms:

\begin{protocolbox}{Tag based aggregate signature protocol ($\ATS$)} {ats}

  \begin{description}

    \item[$(\sk, \pk) \sample \ATSgen(\secparam)$ :] On input the security
      parameter, do the following:

      \begin{enumerate}

        \item Sample $\sk \sample \Zp$.

        \item Return $(\sk, \pk = \ggen^{\sk})$.

      \end{enumerate}

    \item[$\signature \sample \ATSsign(\sk, \Tag, \msg)$ :] On input secret key
      $\sk$ and pair of tag and message $(\Tag, \msg)\in
      \{0,1\}^*\times\{0,1\}^*$, do:

      \begin{enumerate}

        \item Sample $\upsilon \sample \Zp$.

        \item Return $\signature = (\ggen^{\upsilon}, \hashone(\msg)^{\sk} \cdot
          \hashtwo(\Tag)^\upsilon)$.

      \end{enumerate}

    \item[$\signature' \sample \ATSrandomise(\Tag, \signature) $ :] On input
      tag $\Tag$, signature $\signature$, do:

      \begin{enumerate}

        \item Parse $\signature = (\signatureone, \signaturetwo)$.

        \item Sample $\upsilon' \sample \Zp$.

        \item Return $\signature' = (\signatureone \cdot  \ggen^{\upsilon'},
          \signaturetwo \cdot \hashtwo(\Tag)^{\upsilon'})$.

      \end{enumerate}

    \item[$ \{0,1\} \leftarrow \ATSverify(\pk, \Tag, \msg, \signature) $ :] On
      input public key $\pk$, tag $\Tag \in \{0,1\}^*$, message $\msg \in
      \{0,1\}^*$, and signature $\signature$, do:

      \begin{enumerate}

        \item Parse $\signature = (\signatureone, \signaturetwo)$.

        \item Return $e(\signaturetwo,\ggen) \checkcheck e(\hashone(\msg), \pk)
          \cdot e(\hashtwo(\Tag), \signatureone)$.

      \end{enumerate}

    \item[$\signature \leftarrow \ATSaggregate((\signature_i)_{i \in
      \ISSUERS})$:] On input signatures $(\signature_i)_{i \in \ISSUERS}$, for
      set $\ISSUERS \subseteq [1, n]$, do:

      \begin{enumerate}

        \item Parse $\signature_i = (\signatureone_i, \signaturetwo_i)$ for
          each $i \in \ISSUERS$.

        \item Return $\signature = (\prod_{i\in\ISSUERS} \signatureone_i, \
          \prod_{i\in\ISSUERS} \signaturetwo_i)$.

      \end{enumerate}

    \item[$\{0,1\} \leftarrow \ATSverifyagg((\pk_i)_{i\in \ISSUERS},
      \Tag,(\msg_i)_{i\in \ISSUERS},\signature) $ :] On input public keys
      $(\pk_i)_{i\in \ISSUERS}$, tag $\Tag$, messages $(\msg_i)_{i\in
      \ISSUERS}$, and aggregate signature $\signature$, the aggregate
      verification algorithm outputs the result of the following check:

      \begin{enumerate}

        \item Parse $\signature = (\signatureone, \signaturetwo)$.

        \item Return $e(\signaturetwo, \ggen) \checkcheck e (\hashtwo(\Tag),
          \signatureone) \cdot \prod_{i\in\ISSUERS} e(\hashone(\msg_i), \pk_i)$.

      \end{enumerate}

  \end{description} \end{protocolbox}

  \begin{assumptionbox}{Bilinear Computational Diffie-Hellman (BCDH)}{bcdh}

    The BCDH assumption holds if, given $(\gen, \gen^\gamma, \gen^\varepsilon, \gen^\chi, \ggen,
    \ggen^\gamma, \ggen^\varepsilon, \ggen^\chi)$ for uniformly random $\gamma, \varepsilon, \chi \sample \Zp$, no
    PPT adversary can output $e(\gen, \ggen)^{\gamma\varepsilon\chi}$ with non-negligible
    probability.

  \end{assumptionbox}

  \begin{assumptionbox}{Modified Bilinear Computational Diffie-Hellman
    (m-BCDH)}{mbcdh}

    The m-BCDH assumption holds if, given $(\gen, \gen^{1/\varepsilon}, \ggen, \ggen^\gamma,
    \ggen^\varepsilon)$ for uniformly random $\gamma, \varepsilon \sample \Zp$, no PPT adversary can
    output $e(\gen, \ggen)^{\gamma\varepsilon}$ with non-negligible probability.

  \end{assumptionbox}

  This signature is existentially unforgeable under tag-based aggregation in
  the random oracle model under the Bilinear Computational Diffie-Hellman
  (BCDH) assumption and a modified variant of the same assumption. To the best
  of our knowledge, the modified BCDH assumption has not been explicitly
  introduced in prior work. However, it fits within the Uber assumption
  framework and can be shown to hold in the generic bilinear group
  model~\cite{EC:BonBoyGoh05,PAIRING:Boyen08}.  This is shown
  in~\Fref{sec:agg-creds-security}.

  \begin{theorem}\label{thm:agg-creds-thm-1} If the BCDH and m-BCDH
  assumptions hold, then our tag based aggregate signature scheme is
  existentially unforgeable under tag based aggregation in the random oracle
  model.  \end{theorem}

\subsection{Our aggregate anonymous credentials scheme ($\AAC$)}

\paragraph{Overview.} In this instantiation, each user has a pair of secret and
public keys $(\sk, \pk = \gen^\sk)$.  The digital signature $\DS$ uses AGHO
structure preserving signatures~\cite{C:AGHO11}. Every user is assigned a
unique tag $\Tag \in \Zp$. Since structure-preserving signatures signs vectors 
of group elements rather
than scalars, we take $\tid = \hashtwo(\Tag)$ and it is the pair $(\pk, \tid)$
that the registration authority signs, with $\pk$ the user's public key.
Every other computation in the protocol uses $\Tag$ directly. The aggregate
tag-based signature $\ATS$ is instantiated with our construction.

\paragraph{The AGHO structure preserving signature scheme.} The AGHO scheme
supports signing vectors of elements in $\group{1}$, $\group{2}$, or both.  In
the following, we only describe signing elements of $\group{1}$. We use a
structure preserving signature scheme specifically because it has public
keys, messages, and signatures consisting of elements of the same group.
This allows us to build a more efficient protocol; we explain how in more
detail once the registration algorithm has been described.

\begin{primitivebox} {AGHO structure-preserving signatures} {agho}

  \begin{description}

    \item[$ (\sk, \pk) \sample \SPSgen(\secparam, \ell) $ :] On input a
      security parameter $\secparam$ and vector length $\ell$, the generation
      algorithm outputs the following key pair:
         %
      \begin{gather*}
     %
        (\mu_1, \dots, \mu_\ell,  \tau) \sample \Zp^{\ell+1}; \quad \sk = (\mu_1, \dots,
        \mu_\ell, \tau); \\
      %
        \pk = (\widetilde{u}_1, \dots, \widetilde{u}_\ell, v) \quad\text{with}\quad
        \widetilde{u}_i = \ggen^{\mu_i} \in \group{2}, v = \gen^\tau \in \group{1}.
        %
      \end{gather*}
     %

    \item[$\SPSsig \sample \SPSsign(\sk, \vec{m})$ :] On input the secret key
      $\sk = (\mu_1, \dots, \mu_\ell, \tau)$ and message $\vec{m} = (m_1, \dots,
      m_\ell) \in \group{1}^\ell$, the signing algorithm randomly samples $\omega
      \sample \Zp$ and outputs the signature constructed as follows:
    %
      \begin{equation*}
     %
        \signature =  (\widetilde{r}, \widetilde{s}, t) = (\ggen^\omega, \widetilde{r}^\tau,
        (\gen \cdot \prod_{i=1}^\ell m_i^{-\mu_i})^{1/\omega}).
    %
      \end{equation*}
     %

    \item[$\SPSsig' \sample \SPSrandomise(\SPSsig)$ :] On input signature
      $\SPSsig = (\widetilde{r}, \widetilde{s}, t)$, the randomisation algorithm
      randomly samples $\omega' \sample \Zp$ and outputs the signature formed as
      follows:
    %
      \begin{equation*}
        %
        \SPSsig' = (\widetilde{r}', \widetilde{s}', t') = ( \widetilde{r}^{\omega'},
        \widetilde{s}^{\omega'}, t^{\frac{1}{\omega'}}).
    %
      \end{equation*}
     %

    \item[$ \{0,1\} \leftarrow \SPSverify(\pk, \vec{m}, \SPSsig)$ :] On input a
      public key $\pk$, message $\vec{m} = (m_1, \dots, m_\ell) \in
      \group{1}^\ell$, and signature $\SPSsig = (\widetilde{r}, \widetilde{s}, t)$,
      the verification algorithm outputs $1$ if both of the following pairing
      equations hold:
    %
      \begin{equation*}
    %
        e(v, \widetilde{r}) = e(\gen, \widetilde{s}) \quad \wedge \quad e(t,
        \widetilde{r}) \cdot \prod_{i=1}^\ell e(m_i, \widetilde{u}_i) = e(\gen,
        \ggen).
  %
      \end{equation*}
  %
      And otherwise outputs $0$.

  \end{description}

\end{primitivebox}

The AGHO signature scheme is existentially unforgeable under chosen message
attacks under the $q$-Strong Diffie-Hellman ($q$-SDH) assumption~\cite{}.

\begin{assumptionbox}{$q$-Strong Diffie-Hellman ($q$-SDH)}{qsdh}

  The $q$-SDH assumption holds if, given $(\gen, \gen^{x}, \gen^{x^2}, \dots,
  \gen^{x^q}, \ggen, \ggen^{x})$ for uniformly random $x \sample \Zp$, no PPT
  adversary can output a pair $(c, \gen^{1/(x+c)})$ for any $c \in \Zp
  \setminus \{-x\}$ with non-negligible probability.

\end{assumptionbox}

%%%% anon creds protocol

\begin{protocolbox} {Aggregate anonymous credentials protocol} {paac}

  \begin{description}

    \item[$\crs \sample \AACgen(\secparam)$ :] On input security
      parameter $\secparam$, do the following:

      \begin{enumerate}

        \item Sample $(\rask,\rapk) \sample \SPSgen(\secparam,2)$.

        \item For each issuer $i$, sample $(\isk{i},\ipk{i}) \sample
          \ATSgen(\secparam)$.

        \item Return $\crs \leftarrow
          (\gen,\ggen,\hashone,\hashtwo,\rapk,\{\ipk{i}\}_i)$.

      \end{enumerate}

    \item[$(\sk,\regcred) \sample \AACregister(\crs)$ :] On input
      $\crs$ from the user, do:

      \begin{enumerate}

        \item The user samples $\sk \sample \Zp$ and sets $\pk \leftarrow
          \gen^\sk$.

        \item The RA samples a fresh nonce $\nonce \sample \Zp$ and sends it
          to the user, who sends back $\pk$ together with a proof of knowledge
          of $\sk$ on $\nonce$.

        \item If the proof does not verify, the RA aborts.

        \item Otherwise, the RA samples $\Tag \sample \Zp$ uniformly at random,
          makes sure it is unique, sets $\tid \leftarrow \hashtwo(\Tag)$, and
          computes $\psi \sample \SPSsign(\rask,(\pk,\tid))$.

        \item Return $\sk$ to the user, and $\regcred \leftarrow
          (\Tag,\pk,\psi)$ to the user.

      \end{enumerate}

    \item[$\cred \sample \AACissue(\crs,\isk{},\regcred,\att{})$ :] On input
      $(\crs,\isk{i},\regcred,\att{i})$ from the user, with $\atv{i}$ the value
      of attribute $\att{i}$ held by the user, do:

      \begin{enumerate}

        \item Parse $\regcred = (\Tag,\pk,\psi)$ and compute $\tid \leftarrow
          \hashtwo(\Tag)$.

        \item If $\SPSverify(\rapk,(\pk,\tid),\psi) \neq 1$, abort.
          (This instantiates the generic $\DSverify$.)

        \item Otherwise, the user proves knowledge of $\sk$ underlying $\pk$ to
          issuer $i$; if this fails, abort.

        \item Otherwise, compute $\signature_i \sample \ATSsign(\isk{i},\Tag,\atv{i})$.

        \item Return $\cred \leftarrow (\atv{i},\signature_i)$ to the user.

      \end{enumerate}

    \item[$\aggcred \leftarrow \AACaggregate(\crs,\cred_1,\ldots,\cred_n)$ :] On
      input credentials $\cred_1,\ldots,\cred_n$ from the user, all issued
      under the same tag $\Tag$, do:

      \begin{enumerate}

        \item Parse each $\cred_j = (\atv{j},\signature_j)$ for $j \in [1,n]$.

        \item Compute $\signature \leftarrow
          \ATSaggregate(\{\signature_j\}_{j\in[1,n]})$.

        \item Return $\aggcred \leftarrow
          (\Tag,\{\atv{j}\}_{j\in[1,n]},\signature)$.

      \end{enumerate}

  \end{description}

\end{protocolbox}

\begin{protocolbox} {Aggregate anonymous credentials protocol continued}
  {paac-2}

  \begin{description}

    \item[$\proof \sample \AACpresent(\crs,\aggcred,\msg)$ :] On input
      $\aggcred$ and message $\msg$ from the verifier, do:

      \begin{enumerate}

        \item Parse $\aggcred = (\Tag,\{\atv{j}\}_j,\signature)$, and retrieve
          the user's stored $(\sk,\pk,\psi)$.

        \item Randomise $\signature' \sample \ATSrandomise(\Tag,\signature)$
          and $\psi' \sample \SPSrandomise(\psi)$.

        \item Sample $\delta \sample \Zp^*$ and compute the $\delta$-blinded
          values $\bar{\gen}, \bar{v}, \bar{t}, \bar{\signaturetwo},
          \bar{\tid}, \pk', \{H_i\}_{i\in\ISSUERS}$.

        \item Set $\inst = (\pk', \{\atv{i}, H_i\}_{i\in\ISSUERS}, \bar{\gen})$
          and $\witness = (\delta, \sk)$, and compute $\eta \sample
          \NIZKprove(\crs, \inst, \witness, \msg)$, a proof that $\rel(\inst,
          \witness) = 1$ for the relation $\rel$ of~\eqref{eq:proof-rel}.

        \item Return $\proof = (\eta, \bar{\gen}, \bar{v}, \widetilde{r}',
          \widetilde{s}', \bar{t}, \signatureone', \bar{\signaturetwo},
          \bar{\tid}, \pk', \{\atv{i}, H_i\}_{i\in \ISSUERS})$.

      \end{enumerate}

    \item[$\{0,1\} \leftarrow \AACverify(\crs,\msg,\proof)$ :] On input
      $(\msg,\proof)$ from a verifier, parse $\proof = (\eta, \bar{\gen},
      \bar{v}, \widetilde{r}', \widetilde{s}', \bar{t}, \signatureone',
      \bar{\signaturetwo}, \bar{\tid}, \pk', \{\atv{i}, H_i\}_{i\in\ISSUERS})$
      and do:

      \begin{enumerate}

        \item If $\bar{\gen} = 1$, return $0$.

        \item Set $\inst = (\pk', \{\atv{i}, H_i\}_{i\in\ISSUERS}, \bar{\gen})$.
          If $\NIZKverify(\crs, \inst, \eta, \msg) = 0$, return $0$.

        \item Otherwise, if any of the following three pairing checks fails, return
          $0$:
          %
          \begin{gather*}
          %
            e(\bar{\signaturetwo}, \ggen) = \prod_{i\in \ISSUERS}e(H_{i},
            \ipk{i})\cdot e(\bar{\tid}, \signatureone') \ ; \\
            %
            e(\bar{\gen}, \ggen) = e(\bar{t}, \widetilde{r}')\cdot e(\pk',
            \widetilde{u}_1)\cdot e(\bar{\tid}, \widetilde{u}_2) \ ; \\
            %
            e(\bar{v}, \widetilde{r}')= e(\bar{\gen}, \widetilde{s}').
          %
          \end{gather*}
          %

        \item Otherwise, return $1$.

      \end{enumerate}

  \end{description}

\end{protocolbox}

The protocol instantiation is described in full in Protocol~\ref{prot:paac}. The
protocol instantiation works as follows:

\begin{description}

  \item[$\crs \sample \AACgen(\secparam)$ :] The registration authority sets up
    a key pair for the structure preserving signature scheme, which it will
    later use to certify each user's public key and tag.  Messages and
    signatures are vectors of elements of the same group as all other elements
    of this protocol. This allows users to prove knowledge of valid
    registration signatures with an efficient sigma protocol over discrete
    logarithms, rather than needing a general purpose zero knowledge proof.
    Each issuer sets up a key pair for the tag based aggregate signature scheme
    of~\Fref{sec:tag-scheme}.
    
  \item[$(\sk,\regcred) \sample \AACregister(\crs)$ :] The user
    generates a key pair and proves knowledge of the secret key to the
    registration authority. If the proof verifies, the registration authority
    assigns the user a fresh scalar tag $\Tag \sample \Zp$, but signs the
    user's public key together with the group element $\tid = \hashtwo(\Tag)$
    rather than $\Tag$ itself so that we can make use of the structure-preserving 
    signature scheme, allowing us to use sigma protocols as zero knowledge 
    proofs of knowledge, and second, so that we can recover $\Tag$ from $\tid$ 
    during extraction through programming the random oracle in the security proof.
    All other computations in the protocol, including the tag-based aggregate
    signature scheme, use the tag $\Tag$ directly.
    
  \item[$\cred \sample \AACissue(\crs,\isk{},\regcred,\att{})$ :] To request a
    credential, the user first shows the issuer their registration signature
    and proves knowledge of the secret key corresponding to the signed public
    key. If both checks pass, and the issuer is satisfied that the user is
    entitled to the credential, the issuer signs the tag and attribute value.
    
  \item[$\aggcred \leftarrow \AACaggregate(\crs,\cred_1,\ldots,\cred_n)$ :] The
    user aggregates all of the individual signatures they hold under the same
    tag into a single aggregate signature. This step is entirely local.
    
  \item[$\proof \sample \AACpresent(\crs,\aggcred,\msg)$ :] The user
    rerandomises their aggregate signature and registration signature, then
    computes a proof using the sigma protocol enabled by AGHO; the full
    derivation follows below.
    
  \item[$\{0,1\} \leftarrow \AACverify(\crs,\msg,\proof)$ :] Verification of
    the credential presentation corresponds to verification of the zero
    knowledge proofs presented.
    
\end{description}

\begin{theorem}\label{thm:agg-creds-thm-2}

  If the $q$-SDH assumption holds for the AGHO signature scheme, the BCDH and
  m-BCDH assumptions hold for the tag based aggregate signature scheme
  of~\Fref{sec:tag-scheme}, and the DDH assumption holds in $\group{1}$, then
  the aggregate anonymous credentials protocol
  of Protocol~\ref{prot:paac}~and~Protocol~\ref{prot:paac-2} securely realises ideal
  functionality $\F{\AAC}$ in the random oracle model.

\end{theorem}

\begin{assumptionbox}{Bilinear Decisional Diffie-Hellman (BDDH)}{bddh}

  The BDDH assumption holds if no PPT adversary can distinguish $(\gen, \gen^x,
  \gen^y, \ggen, \ggen^z, e(\gen, \ggen)^{xyz})$ from $(\gen, \gen^x, \gen^y,
  \ggen, \ggen^z, \Delta)$ for uniformly random $x, y, z \sample \Zp$ and
  $\Delta \sample \group{T}$.

\end{assumptionbox}


\subsubsection*{Proving knowledge.}\label{sec:proof-inst} The user proves that they possess a valid
registration signature, an aggregate credential under a single tag, and the
secret key corresponding to the registered public key corresponding to the
given tag, without revealing either the tag or the public key.

Let $\rapk = (\widetilde{u}_1, \widetilde{u}_2, v)$ be the registration authority's
AGHO public key and $\ipk{i} = \ggen^{\xi_i}$ issuer $i$'s public key. We set the
following:
%
\begin{equation*}
%
  \crs = (\gen, \ggen), \qquad
  \inst = (\rapk, \{\atv{i}, \ipk{i}\}_{i \in \ISSUERS}), \qquad
  \witness = (\sk, \pk, \tid, \signature, \SPSsig),
%
\end{equation*}
%
where $\pk = \gen^{\sk}$ is the user's public key, $\tid = \hashtwo(\Tag)$,
$\signature = (\signatureone, \signaturetwo)$, and $\SPSsig = (\widetilde{r},
\widetilde{s}, t)$.

To present set of attribute values $\{\atv{i}\}_{i \in\ISSUERS}$ after
receiving $\msg$ from the verifier, the user aggregates the corresponding
tag-based signatures and produces a proof for $\msg$ with respect to relation
$\rel$ defined as $\rel( \inst, \witness) = 1$ if and only if the following
equalities hold:
%
\begin{equation}\label{eq:proof-original}
\begin{aligned}
  e(v, \widetilde{r}) &= e(\gen, \widetilde{s}) \ ; &
  e(t, \widetilde{r}) \cdot e(\gen^{\sk}, \widetilde{u}_1)\cdot e(\tid,  \widetilde{u}_2) &= e(\gen, \ggen) \ ; \\
  \pk &= \gen^{\sk} \ ; &
  e(\signaturetwo, \ggen) &= e(\tid, \signatureone)\cdot\prod_{i\in\ISSUERS}e(\hashone(\atv{i}), \ipk{i}).
\end{aligned}
\end{equation}
%
Since we assume that the
registration authority only issues signatures on well-formed messages, pairs
$(\pk, \hashtwo(\Tag))$ for $\Tag \in \Zp$, we do not require the user to prove
knowledge of the pre image such that $\tid = \hashtwo(\Tag)$. This results in
a more efficient proof.

To reduce the computational cost incurred by this proof, we randomise the
signatures and equations in a way that removes the need for producing zero
knowledge proofs over pairing equations in $\group{T}$, and relies on zero
knowledge proofs of knowledge and equality of discrete logarithms of group
elements in $\group{1}$.

The user first randomises signatures $\signature' \leftarrow
\ATSrandomise(\signature)$ and $\SPSsig' \leftarrow \SPSrandomise(\SPSsig)$.
Let $\signature' = (\signatureone', \signaturetwo')$ and $\SPSsig' =
(\widetilde{r}', \widetilde{s}', t')$.  The user then samples a random $\delta \in
\Zp^*$ and computes the following:
%
\begin{gather*}
%
  \forall i \in \ISSUERS: \quad H_i = \hashone(\atv{i})^{\delta} \quad ;
  \quad \bar{\gen} = \gen^{\delta} \quad ;  \quad \bar{t} = {t'}^{\delta} \quad
  ;  \quad \bar{\signaturetwo} = {\signaturetwo'}^{\delta} \ ; \\
%
  \bar{\tid} = \tid^{\delta} \quad ; \quad \bar{v} = v^{\delta} \quad ;
  \quad \pk' = \pk^{\delta}.
%
\end{gather*}

We redefine $\crs = \gen$, $\inst = (\pk', \{\atv{i}, H_i\}_{i\in \ISSUERS},
\bar{\gen})$, and $\witness = (\delta, \sk)$, and let $\rel$ be the relation
defined as $\rel(\inst, \witness) = 1$ if and only if the following holds:
%
\begin{equation} \label{eq:proof-rel}
%
  \bar{\gen} = \gen^{\delta} \quad ; \quad \bar{v} = v^{\delta} \quad ;  \quad
  H_i = \hashone(\atv{i})^{\delta} \quad ; \quad \pk' = \bar{\gen}^{\sk}.
%
\end{equation}
 %
Next, the user produces $\eta$, a proof for $\rel$, and defines 
the presentation proof as follows:
%
\begin{equation} \label{eq:proof-phi}
%
  \proof = (\eta, \bar{\gen}, \bar{v}, \widetilde{r}', \widetilde{s}', \bar{t},
  \signatureone', \bar{\signaturetwo}, \bar{\tid}, \pk', \{\atv{i},
  H_i\}_{i\in \ISSUERS}).
%
\end{equation}
 %
The verifier accepts $\proof$ if and only if $\bar{\gen} \neq 1$, 
$\eta$ verifies, and the
following equalities hold:
%
\begin{gather}
 %
  \label{eq:proof-verify}
  e(\bar{\signaturetwo}, \ggen)  = \prod_{i\in \ISSUERS}e(H_{i},
  \ipk{i})\cdot e(\bar{\tid}, \signatureone') \ ; \ \\
  %
  e(\bar{\gen}, \ggen) = e(\bar{t}, \widetilde{r}')\cdot e(\pk',
  \widetilde{u}_1)\cdot e(\bar{\tid}, \widetilde{u}_2) \ ; \\
  %
  e(\bar{v}, \widetilde{r}')= e(\bar{\gen}, \widetilde{s}').
%
\end{gather}
%

Next we describe how $\eta$ is computed.  This is followed by showing that if
the user is able to produce a valid $\eta$, then we can extract a valid witness
for the relation defined in~\fref{eq:proof-original}.  We prove that under the
DDH assumption in $\group{1}$, $\proof$ is simulatable, it does not leak any
information about the witness.

\begin{lemma}\label{lem:agg-creds-extract}

  If a user produces a presentation  that verifies, then a witness for
  relation~\eqref{eq:proof-original} can be extracted.

\end{lemma}

\begin{lemma}\label{lem:agg-creds-simulate}

  If the DDH assumption holds in $\group{1}$, then a valid presentation can be
  simulated without access to a witness.

\end{lemma}

\paragraph{$\eta$ Computation.}\label{sec:agg-creds-eta} This section gives
the concrete construction that $\NIZKprove$ and $\NIZKverify$ instantiate
in Protocol~\ref{prot:paac-2}. We first describe the
interactive zero-knowledge proof  underlying the computation of $\eta$, and
then show how to apply the Fiat-Shamir heuristics to the proof to finally
obtain $\eta$.

The user proves knowledge of a valid witness $\witness = (\delta, \sk)$ for
the relation depicted in equation~\eqref{eq:proof-rel}, by executing the
following:

\begin{itemize}

  \item Select two random numbers $\alpha$ and $\beta$ in $\Zp$;

  \item Compute and send the following tuple to the verifier:
  %
    \begin{equation*}
  %
      (f, h, \{k_i\}_{i \in \ISSUERS}, w) = (\gen^{\alpha}, v^{\alpha},
      \{\hashone(\atv{i})^{\alpha}\}_{i \in \ISSUERS}, \bar{\gen}^{\beta});
    %
    \end{equation*}
    %

  \item On receiving challenge $c$ from the verifier, compute and send $\pi_1 =
    \alpha + c\cdot\delta$ and $\pi_2 = \beta + c\cdot\sk$.

\end{itemize}

The verifier accepts $\pi_1$ and $\pi_2$ if the following equalities hold:
%
\begin{equation*}
%
  \gen^{\pi_1} = f\cdot\bar{\gen}^c, \qquad v^{\pi_1} = h\cdot\bar{v}^c, \qquad
  \hashone(\atv{i})^{\pi_1} = k_i\cdot H_i^c, \qquad \bar{\gen}^{\pi_2} =
  w\cdot\pk'^c.
%
\end{equation*}
%
To produce $\eta$ on input of a message $\msg$, we apply the Fiat-Shamir
heuristic to generate the challenge $c$.  Let $\hash: \{0,1\}^* \rightarrow
\Zp$ be a cryptographic hash function. The user computes $\eta$ by first
computing tuple $(f, h, \{k_i\}_{i \in \ISSUERS}, w)$, followed by $c =
\hash(f, h, \{k_i\}_{i \in \ISSUERS}, w, \msg)$, and outputting $\eta = (f,
h, \{k_i\}_{i \in \ISSUERS}, w, \pi_1, \pi_2)$.

\subsubsection*{Extensions.} The construction supports also several
straightforward extensions.

New issuers can join the system without impacting the construction.  In fact,
the aggregation of credentials and their presentation is agnostic to the number
of issuers.

The current construction assumes that each issuer produces a credential for a
single attribute type.  The current construction captures settings where
issuers generate credentials for multiple attribute types by having each issuer
hold different public key for each attribute type.  Alternatively, if it is
desirable to have the issuer hold a single public key regardless of the number
of attribute types it is responsible for, then we modify the way the issuer
signs. Instead of signing pair $(\Tag, \atv{})$, the issuer signs $(\Tag,
\att{} \| \atv{})$, with $\att{}$ the attribute type and $\atv{}$ the attribute
value.

The construction can be extended to support epoch-based revocation, in the
following way.  Instead of signing pair $(\pk, \tid)$, the registration
authority signs triple $(\pk, \tid, \gen^{\epoch})$, where $\epoch$ identifies
the current epoch.  Issuers, on
their end, sign pair $(\Tag, \epoch \| \att{} \| \atv{})$, instead of pair
$(\Tag, \att{} \| \atv{})$; this ties the attribute credential to epoch
$\epoch$.  During credential presentation, users prove that they hold a valid
signature from the registration authority and valid aggregate tag-based
signatures from the relevant issuers, for the current epoch.

We leave the study of accumulator-based revocation and the revocation of
individual attribute credentials to future work.